\documentclass[aspectratio=169,12pt]{beamer}
\usetheme{Madrid}
\usecolortheme{dolphin}
\usefonttheme{professionalfonts}
\setbeamertemplate{navigation symbols}{}
\setbeamertemplate{footline}{}

\usepackage[T1]{fontenc}
\usepackage[utf8]{inputenc}
\usepackage{graphicx}
\usepackage{amsmath,amssymb}
\usepackage{hyperref}
\usepackage{siunitx}
\usepackage{physics}
\usepackage{cancel}
\graphicspath{{figures/}}

\title[Chapter 6]{Chapitre 6: Systèmes à 2 qubits et intrication}
\author{JAMOTTE Maxime, SCHOONEN Cédric}
\institute{Digital Learning Hub}
\date{}

\begin{document}

\begin{frame}
  \titlepage
\end{frame}

% \begin{frame}{Table des matières}
%   \tableofcontents
% \end{frame}

% \section{}

\begin{frame}{Objectifs du chapitre 6}
  \begin{enumerate}[<+->]
    \item Représentation tensorielle d'un état à 2 qubits.
    \\~\\
    \item Portes quantiques sur 2 qubits
    \\~\\
    \item Générer une paire de Bell
    \\~\\
    \item Conséquences de l'intrication sur la mesure
  \end{enumerate}
\end{frame}

% ============================================================
\section{Notation tensorielle}
% ============================================================

\begin{frame}{Représentation tensorielle d'un état à 2 qubits}
  \begin{itemize}[<+->]
    \item Deux qubits dans l'état "0": $\ket{0}\!\otimes\!\ket{0}$
    \\~\\
    \item Un qubit dans l'état "0" et l'autre dans "1": $\ket{0}\!\otimes\!\ket{1}$
    \\~\\
    \item Superposition de ces deux états: $\ket{0}\!\otimes\!\ket{0} + \ket{0}\!\otimes\!\ket{1}$.
  \end{itemize}
\end{frame}

\begin{frame}{Propriétés du produit tensoriel}
  \begin{itemize}[<+->]
    \setlength{\itemsep}{0.6em}
    \item \textbf{Distributivité} sur l'addition:
    $$\ket{0}\!\otimes\!\big(\alpha\ket{0}+\beta\ket{1}\big)
      = \alpha\!\ket{0}\!\otimes\!\ket{0}+\beta\!\ket{0}\!\otimes\!\ket{1}$$
    \item \textbf{Associativité}: pour 3 qubits ou plus, l'ordre de regroupement n'importe pas:
    $$\big(\ket{a}\!\otimes\!\ket{b}\big)\!\otimes\!\ket{c}
      = \ket{a}\!\otimes\!\big(\ket{b}\!\otimes\!\ket{c}\big)
      = \ket{abc}$$
    \item Ces règles permettent de développer et simplifier des états multi-qubits.
    \item \textbf{Attention}: le produit tensoriel n'est \textbf{pas commutatif}:
     $\ket{0}\!\otimes\!\ket{1} \neq \ket{1}\!\otimes\!\ket{0}$.
  \end{itemize}
\end{frame}

% ============================================================
\section{Base computationnelle et numérotation binaire}
% ============================================================

\begin{frame}{Un grand nombre d'états}
  \begin{itemize}
    \item État général: $\ket{\Psi_2} = a_{00}\ket{00}+a_{01}\ket{01}+a_{10}\ket{10}+a_{11}\ket{11}$
    \\~\\
    \item Numérotation binaire des états de la base computationnelle: $00, 01, 10, 11$
    \\~\\
    \item Le nombre d'états encodables grandit exponentiellement vite
    \\~\\
    \item Pour $n$ qubits: $2^n$ états (1, 4, 8, 16, 32, $\cdots$) $\rightarrow$ environ un million pour 20 qubits!.
  \end{itemize}
\end{frame}

% ============================================================
\section{Représentation vectorielle}
% ============================================================

\begin{frame}{Représentation vectorielle}
  \begin{itemize}
    \item État général: $\ket{\Psi_2} = a_{00}\ket{00}+a_{01}\ket{01}+a_{10}\ket{10}+a_{11}\ket{11}$
  \end{itemize}
  ~\\~\\
  $\qquad \ket{\Psi_2} = \begin{pmatrix} 
    a_{00} \\
    a_{01} \\
    a_{10} \\
    a_{11} \\
\end{pmatrix} \longrightarrow \ket{00} = \begin{pmatrix} 
    1\\
    0\\
    0\\
    0\\
\end{pmatrix} , \ket{01} = \begin{pmatrix} 
    0\\
    1\\
    0\\
    0\\
\end{pmatrix}, \ket{10} = \begin{pmatrix} 
    0\\
    0\\
    1\\
    0\\
\end{pmatrix} , \ket{11} = \begin{pmatrix} 
    0\\
    0\\
    0\\
    1\\
\end{pmatrix}  $
\end{frame}

% ============================================================
\section{Produits tensoriels de matrices}
% ============================================================

\begin{frame}{Action des opérateurs sur un état à 2 qubits}
  \begin{itemize}[<+->]
    \item On peut appliquer des opérateurs séparément sur chaque qubit.
    \\~\\
    \item L'action séparée de deux matrices $M_2$ et $M_1$ sur chaque qubit se traduit par un produit tensoriel de matrices:
    \[
      (M_2\ket{\psi_2})\!\otimes\!(M_1\ket{\phi_1}) = (M_2\!\otimes\! M_1)(\ket{\psi_2}\!\otimes\!\ket{\phi_1})
    \]
  \end{itemize}
\end{frame}

\begin{frame}{Hadamard en parallèle}
  \begin{itemize}[<+->]
    \setlength{\itemsep}{0.55em}
    \item Pour appliquer $H$ à chaque qubit indépendamment, on utilise $H\otimes H$:
    $$(H\otimes H)\ket{00} = H\ket{0}\otimes H\ket{0}
     = \frac{\ket{0}+\ket{1}}{\sqrt{2}}\otimes\frac{\ket{0}+\ket{1}}{\sqrt{2}}$$
    $$= \frac{1}{2}\big(\ket{00}+\ket{01}+\ket{10}+\ket{11}\big)$$
    \item Superposition \textbf{uniforme} sur les 4 états de base.
  \end{itemize}
\end{frame}

\begin{frame}{Hadamard en parallèle}
  \begin{equation*}
    \frac{1}{\sqrt{2}}
    \begin{pmatrix}
      1 &  1 \\
      1 & -1 \\
    \end{pmatrix}
    \otimes
    \frac{1}{\sqrt{2}}
    \begin{pmatrix}
      1 &  1 \\
      1 & -1 \\
    \end{pmatrix}
    % \begin{pmatrix}
    %   1 \\
    %   0 \\
    %   0 \\
    %   0 \\
    % \end{pmatrix}
    = 
    \frac{1}{2}
    \begin{pmatrix}
      1 &  1 &  1 &  1 \\
      1 & -1 &  1 & -1 \\
      1 &  1 & -1 & -1 \\
      1 & -1 & -1 &  1 \\
    \end{pmatrix}
    % \begin{pmatrix}
    %   1 \\
    %   0 \\
    %   0 \\
    %   0 \\
    % \end{pmatrix}
    % =
    % \frac{1}{2}
    % \begin{pmatrix}
    %   1 \\
    %   1 \\
    %   1 \\
    %   1 \\
    % \end{pmatrix}
  \end{equation*}
\end{frame}

\begin{frame}{Opérateurs non-décomposables en produit tensoriel}
  \begin{itemize}[<+->]
    \setlength{\itemsep}{0.6em}
    \item Certaines opérations à 2 qubits ne s'écrivent \textbf{pas} comme $M_2 \otimes M_1$.
    \\~\\
    \item Trois opérateurs importants : \textbf{SWAP}, \textbf{CX} (CNOT) et \textbf{CZ}.
  \end{itemize}
\end{frame}

\begin{frame}{Porte SWAP}
  \begin{itemize}[<+->]
    \setlength{\itemsep}{0.55em}
    \item Permute les deux qubits: $\text{SWAP}\ket{ab} = \ket{ba}$.\\~\\
    \item Matrice dans la base computationnelle $\{\ket{00},\ket{01},\ket{10},\ket{11}\}$:
    $$\text{SWAP} = \begin{pmatrix} 1&0&0&0 \\ 0&0&1&0 \\ 0&1&0&0 \\ 0&0&0&1 \end{pmatrix}$$
  \end{itemize}
\end{frame}

\begin{frame}{Porte CNOT (controlled-X)}
    \begin{itemize}[<+->]
      \setlength{\itemsep}{0.55em}
      \item Inverse le qubit \textbf{cible} si le qubit \textbf{contrôle} est $\ket{1}$.
      \item CNOT$\ket{00}=\ket{00}$, CNOT$\ket{01}=\ket{01}$,
       CNOT$\ket{10}=\ket{11}$, CNOT$\ket{11}=\ket{10}$.
      \item Matrice:
      $$\text{CNOT} = \begin{pmatrix} 1&0&0&0 \\ 0&1&0&0 \\ 0&0&0&1 \\ 0&0&1&0 \end{pmatrix}$$
    \end{itemize}
\end{frame}

\begin{frame}{Porte CZ (controlled-Z)}
  \begin{itemize}[<+->]
    \setlength{\itemsep}{0.55em}
    \item Applique $Z$ au qubit cible si le qubit contrôle est $\ket{1}$:
     CZ$\ket{11} = -\ket{11}$.\\~\\
    \item Matrice:
    $$\text{CZ} = \begin{pmatrix} 1&0&0&0 \\ 0&1&0&0 \\ 0&0&1&0 \\ 0&0&0&-1 \end{pmatrix}$$
    \item La porte CZ est \textbf{symétrique}: peu importe quel qubit est le contrôle.
  \end{itemize}
\end{frame}

\begin{frame}
  \begin{itemize}[<+->]
    \item Les opérateurs $C_X^{}$ et $C_Z^{}$ peuvent \textbf{corréler} les qubits entre eux et peuvent créer de l'\textbf{intrication}.\\~\\
    \item La mesure d'un qubit intriqué influence l'autre, même à distance! Voyons cela avec une paire de qubits intriqués.
     \\~\\
  \end{itemize}
\end{frame}

% \section{États à deux qubits}

% \begin{frame}{1) Représentation tensorielle d'un état à 2 qubits}
%   \begin{itemize}[<+->]
%     \item État général $\ket{\Psi_2} = a_{00}\ket{00}+a_{01}\ket{01}+a_{10}\ket{10}+a_{11}\ket{11}$ (vecteur de taille $2^2$).
%     \\~\\
%     \item Produit tensoriel de kets: $\ket{0}_2\!\otimes\!\ket{0}_1=\ket{00}$.
%     \\~\\
%     \item Nombre d'états de base pour $N$ qubits: $2^N$ (ex.: 2 qubits $\rightarrow$ 4 composantes).
%   \end{itemize}
% \end{frame}

% \begin{frame}{1) Représentation tensorielle d'un état à 2 qubits}
%   \begin{itemize}
%     \item État général $\ket{\Psi_2} = a_{00}\ket{00}+a_{01}\ket{01}+a_{10}\ket{10}+a_{11}\ket{11}$ (vecteur de taille $2^2$).
%     \\~\\
%     \item Produit tensoriel de kets: $\ket{0}_2\!\otimes\!\ket{0}_1=\ket{00}$.
%     \\~\\
%     \item Nombre d'états de base pour $N$ qubits: $2^N$ (ex.: 2 qubits $\rightarrow$ 4 composantes).
%   \end{itemize}
%   ~\\~\\
%   $\qquad \qquad\begin{pmatrix} 
%     a_{00}^{} \\
%     a_{01}^{} \\
%     a_{10}^{}\\
%     a_{11}^{} \\
% \end{pmatrix} \longrightarrow \ket{00} = \begin{pmatrix} 
%     1\\
%     0\\
%     0\\
%     0\\
% \end{pmatrix} , \ket{01} = \begin{pmatrix} 
%     0\\
%     1\\
%     0\\
%     0\\
% \end{pmatrix}, \ket{10} = \begin{pmatrix} 
%     0\\
%     0\\
%     1\\
%     0\\
% \end{pmatrix} , \ket{11} = \begin{pmatrix} 
%     0\\
%     0\\
%     0\\
%     1\\
% \end{pmatrix}  $
% \end{frame}

% \begin{frame}{2) États-produits vs intriqués}
%   \begin{itemize}[<+->]
%     \item \textbf{État produit}: $\ket{\Psi_2}=\ket{\phi_2}\!\otimes\!\ket{\psi_1}$; évolution \textbf{indépendante} des qubits dans le circuit. Etat produit = état séparable.
%     \\~\\
%     \item \textbf{Non séparable} = intriqué: mesurer un qubit influence l'autre.
%     \\~\\
%     \item Paires de Bell: (anti)corrélation totale entre les mesures des deux qubits!
%     \begin{align*}
%       \ket{\Phi_{\pm}}=\frac{\ket{00}\pm\ket{11}}{\sqrt{2}}
%       \qquad\qquad
%       \ket{\Psi_{\pm}}=\frac{\ket{01}\pm\ket{10}}{\sqrt{2}}
%     \end{align*}
%   \end{itemize}
% \end{frame}



% \section{Opérations sur plusieurs qubits}

% \begin{frame}{3) Portes quantiques sur deux qubits: portes locales en parallèle}
%   \begin{itemize}[<+->]
%     \item Action séparée: $(M_2\ket{\psi_2})\!\otimes\!(M_1\ket{\phi_1}) = (M_2\!\otimes\! M_1)(\ket{\psi_2}\!\otimes\!\ket{\phi_1})$.
%     \\~\\
%     \item Circuit exemple: deux portes de Hadamard appliquées sur deux qubits initiaux $\ket{0}$ et $\ket{1}$.
%   \end{itemize}
%   \begin{center}
%     \includegraphics[width=0.3\textwidth]{parallel_hadamard.png}
%   \end{center}
% \end{frame}

% \begin{frame}{3) Portes quantiques à deux qubits: porte SWAP}
%   \begin{itemize}[<+->]
%     \item Échange les qubits: $\ket{01}\leftrightarrow\ket{10}$ via une matrice $4\times4$ de permutation.
%     \\~\\
%     \item Ne crée pas d'intrication (seule permutation des composantes).\\~\\
%   \end{itemize}
%   \[S = \begin{pmatrix}
%       1 & 0 & 0 & 0\\
%       0 & 0 & 1 & 0\\
%       0 & 1 & 0 & 0\\
%       0 & 0 & 0 & 1
%     \end{pmatrix}
%   \]
%   \begin{center}
%     \includegraphics[width=0.3\textwidth]{swap_gate.png}
%   \end{center}
% \end{frame}

% \begin{frame}{3) Portes quantiques à deux qubits: porte CNOT (CX)}
%   \begin{itemize}[<+->]
%     \item Contrôle sur le deuxième qubit, cible sur le premier.
%     \item Retourne la cible si le contrôle vaut $\ket{1}$; matrice $4\times4$ non factorisable $\Rightarrow$ peut intriquer.
%   \end{itemize}
%   \[
%     C_X = \begin{pmatrix}
%       1 & 0 & 0 & 0\\
%       0 & 1 & 0 & 0\\
%       0 & 0 & 0 & 1\\
%       0 & 0 & 1 & 0
%     \end{pmatrix}
%   \]
%   \begin{center}
%     \includegraphics[width=0.3\textwidth]{cnot_gate.png}
%   \end{center}
% \end{frame}

% \begin{frame}{3) Portes quantiques à deux qubits: porte CZ}
%   \begin{itemize}[<+->]
%     \item Change le signe du coefficient devant $\ket{11}$ seulement.
%     \\~\\
%     \item Utilisée dans plusieurs algorithmes; peut aussi générer de l'intrication.
%   \end{itemize}~\\
%   \[
%     C_Z = \begin{pmatrix}
%       1 & 0 & 0 & 0\\
%       0 & 1 & 0 & 0\\
%       0 & 0 & 1 & 0\\
%       0 & 0 & 0 & -1
%     \end{pmatrix}
%   \]
%   \begin{center}
%     \includegraphics[width=0.32\textwidth]{cz_gate.png}
%   \end{center}
% \end{frame}

\begin{frame}{États-produits vs intriqués}
  \begin{itemize}[<+->]
    \item \textbf{État produit}: $\ket{\Psi_2}=\ket{\phi_2}\!\otimes\!\ket{\psi_1}$; évolution \textbf{indépendante} des qubits dans le circuit. Etat produit = état séparable.
    \\~\\
    \item \textbf{Non séparable} = intriqué: mesurer un qubit influence l'autre.
    \\~\\
    \item Paires de Bell: (anti)corrélation totale entre les mesures des deux qubits!
    \begin{align*}
      \ket{\Phi_{\pm}}=\frac{\ket{00}\pm\ket{11}}{\sqrt{2}}
      \qquad\qquad
      \ket{\Psi_{\pm}}=\frac{\ket{01}\pm\ket{10}}{\sqrt{2}}
    \end{align*}
    \item Mesure d'une paire Bell $\ket{\Phi_+}$: résultats toujours corrélés ($\ket{00}$ ou $\ket{11}$), chaque issue avec probabilité $0.5$.
  \end{itemize}
\end{frame}

\begin{frame}{Génération des paires de Bell}
    \begin{itemize}[<+->]
      \setlength{\itemsep}{0.55em}
      \item Appliquer $H$ sur le premier qubit puis CNOT:
      $$\ket{00} \xrightarrow{H\otimes I} \frac{\ket{0}+\ket{1}}{\sqrt{2}}\ket{0}$$
      $$\xrightarrow{\text{CNOT}} \frac{\ket{00}+\ket{11}}{\sqrt{2}} = \ket{\Phi^+}$$
      \item C'est un état \textbf{intriqué}: la paire de Bell $\ket{\Phi^+}$.
    \end{itemize}
\end{frame}

% ============================================================
\section{Décomposition de Schmidt}
% ============================================================

\begin{frame}{Tester l'intrication}
  \begin{itemize}[<+->]
    \item L'état $\frac{\ket{00} + \ket{01}}{\sqrt{2}}$ est-il intriqué? 
    \item Non: on peut le réécrire comme un état produit
    \begin{equation*}
      \frac{\ket{00} + \ket{01}}{\sqrt{2}} = \ket{0} \otimes \frac{\ket{0}+\ket{1}}{2}
    \end{equation*} 
    \item Si l'état est intriqué, on obtient toujours plusieurs termes:
    \begin{equation*}
      \frac{\ket{00} + \ket{11}}{\sqrt{2}} = \frac{1}{\sqrt{2}} \ket{0} \otimes \ket{0} + \frac{1}{\sqrt{2}} \ket{1} \otimes \ket{1}
    \end{equation*}
    \item La \textbf{décomposition de Schmidt} donne une procédure pour tester l'intrication.
  \end{itemize}
\end{frame}

\begin{frame}{Mesures partielles (qubit d'Alice)}
  \begin{itemize}[<+->]
    \item Exemple plus intéressant:
    \begin{equation*}
      \ket{\psi_{AB}} = \frac{2\ket{00} + \ket{01} + \ket{10} + \ket{11}}{\sqrt{7}}
    \end{equation*}
    \item Décomposition mettant en évidence le qubit d'Alice:
    \begin{align*}
      \ket{\psi_{AB}} 
      &= \frac{1}{\sqrt{7}} (2\ket{00} + \ket{01}) + \frac{1}{\sqrt{7}} (\ket{10} + \ket{11})\\
      &= \sqrt{\frac{5}{7}} \left( \ket{0} \otimes \frac{2\ket{0}+\ket{1}}{\sqrt{5}} \right) + \sqrt{\frac{2}{7}} \left( \ket{1} \otimes \frac{\ket{0}+\ket{1}}{\sqrt{2}} \right)
    \end{align*}
    \item Les probabilités de mesure d'Alice sont:
    \begin{equation*}
      P(\ket{0}_A) = \frac{5}{7} \qquad\qquad P(\ket{1}_A) = \frac{2}{7}
    \end{equation*}
  \end{itemize}
\end{frame}

\begin{frame}{Mesures partielles (qubit de Bob)}
  \begin{itemize}
    \item Exemple plus intéressant:
    \begin{equation*}
      \ket{\psi_{AB}} = \frac{2\ket{00} + \ket{01} + \ket{10} + \ket{11}}{\sqrt{7}}
    \end{equation*}
    \item Décomposition mettant en évidence le qubit de Bob:
    \begin{align*}
      \ket{\psi_{AB}} 
      &= \frac{1}{\sqrt{7}} (2\ket{00} + \ket{10}) + \frac{1}{\sqrt{7}} (\ket{01} + \ket{11})\\
      &= \sqrt{\frac{5}{7}} \left( \frac{2\ket{0}+\ket{1}}{\sqrt{5}} \otimes \ket{0} \right) + \sqrt{\frac{2}{7}} \left( \frac{\ket{0}+\ket{1}}{\sqrt{2}} \otimes \ket{1} \right)
    \end{align*}
    \item Les probabilités de mesure de Bob sont aussi:
    \begin{equation*}
      P(\ket{0}_B) = \frac{5}{7} \qquad\qquad P(\ket{1}_B) = \frac{2}{7}
    \end{equation*}
  \end{itemize}
\end{frame}

\begin{frame}{Conséquence de l'intrication sur les mesures partielles}
  \begin{itemize}[<+->]
    \item Que se passe-t-il du point de vue d'Alice si Bob mesure $\ket{1}$?
    \begin{align*}
      \ket{\psi_{AB}} 
      &= \sqrt{\frac{5}{7}} \cancel{\left( \frac{2\ket{0}+\ket{1}}{\sqrt{5}} \otimes \ket{0} \right)} + \sqrt{\frac{2}{7}} \left( \frac{\ket{0}+\ket{1}}{\sqrt{2}} \otimes \ket{1} \right)
    \end{align*}
    \item L'état quantique devient
    \begin{align*}
      \ket{\psi_{AB}} 
      &= \frac{\ket{0}+\ket{1}}{\sqrt{2}} \otimes \ket{1} = \frac{1}{\sqrt{2}} \ket{0} \otimes \ket{1} + \frac{1}{\sqrt{2}} \ket{1} \otimes \ket{1}
    \end{align*}
    \item Et les probabilités de mesure d'Alice sont maintenant:
    \begin{equation*}
      P(\ket{0}_A) = \frac{1}{2} \qquad\qquad P(\ket{1}_A) = \frac{1}{2}
    \end{equation*}
  \end{itemize}
\end{frame}

\begin{frame}{Fin du module A4}
    \centering
  \vfill
  Vers les circuits et portes logiques!
  \vfill
\end{frame}

\end{document}
